\section{Grassmann 簇与 Plücker 映射}\label{sec:grassmannians-and-plucker-maps}

Grassmann 簇以 Hermann Grassmann (1809--1877) 命名, 是一族参数化一个固定向量空间的线性子空间的射影簇. 等价地, 它们参数化一个固定射影空间的线性子簇. Grassmann 簇在代数几何的许多专题中起核心作用. 本节的目标是抽象地引入 Grassmann 簇, 描述如何经由 Plücker 映射将它们实现为射影簇, 并略述如何用它们回答有趣的几何问题. 先从集合论定义开始.

\begin{Definition}{作为线性子空间的集合的 Grassmann 簇}\leavevmode\label{def:grassmannian-as-set-of-vector-subspaces}
设 $k,n$ 为满足 $k\le n$ 的非负整数. \emph{Grassmann 簇} $\mathrm{G}(k,n)$ 是向量空间 $K^{n}$ 的所有 $k$ 维线性子空间的集合.
\end{Definition}

例如, 对任意 $n$, $\mathrm{G}(0,n)$ 只是单个点, 因为 $K^{n}$ 只有一个零维线性子空间. 而 $K^{n}$ 中每条过原点的直线对应于 $\mathrm{G}(1,n)$ 中的一个点, 因此 $\mathrm{G}(1,n)$ 与 $\mathbb{P}^{n-1}$ 有自然双射 (Corollary~\ref{prop:projective-space-affine-plus-infinity}).

另一种常见的观点是用射影空间的线性子簇看待 Grassmann 簇, 现在就来定义这种子簇.

\begin{Definition}{射影空间的线性子簇}\leavevmode\label{def:linear-subvariety-of-projective-space}
若存在 $(k+1)$ 维线性子空间 $V\subseteq K^{n+1}$, 使得
\[
  L=\mathbb{P}(V)=\{[a_{0}:\cdots:a_{n}]\mid(a_{0},\ldots,a_{n})\in V\setminus\{0\}\},
\]
就称子集 $L\subseteq\mathbb{P}^{n}$ 为 $\mathbb{P}^{n}$ 的\emph{$k$ 维线性子簇}.
\end{Definition}

一维和二维的线性子簇分别称为\emph{直线}和平面, 余维为一的线性子簇称为\emph{超平面}. 注意, 线性子簇上的仿射锥恢复了 $\mathbb{A}^{n+1}=K^{n+1}$ 中原来给出它的线性子空间. 因此, $\mathbb{P}^{n}$ 的 $k$ 维线性子簇与 $K^{n+1}$ 的 $(k+1)$ 维线性子空间之间有自然双射 (Exercise~\ref{ex:linear-projective-subvarieties-correspond-to-vector-subspaces}). 于是得到 Grassmann 簇的如下另一种解释.

\begin{Definition}{作为线性子簇的集合的 Grassmann 簇}\leavevmode\label{def:grassmannian-as-set-of-projective-subvarieties}
设 $k,n$ 为满足 $k\le n$ 的正整数. \emph{Grassmann 簇} $\mathrm{G}(k,n)$ 是 $\mathbb{P}^{n-1}$ 的所有 $(k-1)$ 维线性子簇的集合.
\end{Definition}

Grassmann 簇为研究射影空间中有关线性子簇的几何问题提供了自然框架. 例如, 其中一个最简单的问题是: 给定 $\mathbb{P}^{3}$ 中的四条直线, 是否存在一条直线与它们都相交? 正如本节末尾将看到的, 答案是\textquotedbl{}是\textquotedbl{}. 而且, 在对给定四条直线作一般性假设的条件下, 恰有两条不同的直线与它们都相交. 研究这一问题及许多相关问题的关键在于深化对 $\mathrm{G}(k,n)$ 的理解, 不再仅把它看成抽象集合, 而是赋予它射影簇结构.

为了把 Grassmann 簇实现为射影簇, 我们利用线性子空间的矩阵表示. 更具体地说, 给定一个 $k$ 维线性子空间 $V\subseteq K^{n}$, 换句话说, 给定一个元素 $V\in\mathrm{G}(k,n)$, 选取 $V$ 的一组基 $v_{1},\ldots,v_{k}$, 并把 $V$ 关于该基的\emph{矩阵表示}定义为以 $v_{1},\ldots,v_{k}$ 为各行的 $k\times n$ 矩阵 $M(v_{1},\ldots,v_{k})$. $V$ 的不同基给出不同的矩阵表示, 但这些矩阵都具有共同的性质, 如下一例所示.

\begin{Example}{表示同一个子空间的不同矩阵}\leavevmode\label{ex:matrix-representations-of-same-plane}
令 $V\subseteq K^{4}$ 为 $v_{1}=(2,1,-3,0)$ 和 $v_{2}=(1,0,1,1)$ 张成的空间. 由于 $v_{1}$ 和 $v_{2}$ 线性无关, 它们构成 $V$ 的一组基, 并且
\[
  M(v_{1},v_{2})=\begin{pmatrix}
    2&1&-3&0\\
    1&0&1&1
  \end{pmatrix}.
\]
然而, 选取 $V$ 的另一组基, 例如 $w_{1}=2v_{1}-v_{2}=(3,2,-7,-1)$ 和 $w_{2}=v_{1}+v_{2}=(3,1,-2,1)$, 得到不同的矩阵表示:
\[
  M(w_{1},w_{2})=\begin{pmatrix}
    3&2&-7&-1\\
    3&1&-2&1
  \end{pmatrix}.
\]
这两个矩阵表示当然不是毫不相干的: 它们由 $v_{1},v_{2}$ 与 $w_{1},w_{2}$ 之间的换基矩阵联系起来:
\[
  \begin{pmatrix}3&2&-7&-1\\3&1&-2&1\end{pmatrix}
  =\begin{pmatrix}2&-1\\1&1\end{pmatrix}
   \begin{pmatrix}2&1&-3&0\\1&0&1&1\end{pmatrix}.
\]
特别是, 从每个矩阵表示中挑出两列, 例如第一列和第四列, 由行列式的乘法性可知
\[
  \det\begin{pmatrix}3&-1\\3&1\end{pmatrix}
  =\det\begin{pmatrix}2&-1\\1&1\end{pmatrix}
   \cdot\det\begin{pmatrix}2&0\\1&1\end{pmatrix}.
\]
因此, $M(w_{1},w_{2})$ 的每个 $2\times2$ 子式都是 $M(v_{1},v_{2})$ 对应的 $2\times2$ 子式乘以 $3$ 得到的. 换句话说, 这两个矩阵各自的六个 $2\times2$ 子式给出 $\mathbb{P}^{5}$ 中的同一个点, 而且 $V$ 的任意其他矩阵表示的 $2\times2$ 子式仍确定 $\mathbb{P}^{5}$ 中的同一个点. 正如下面将正式定义的, $\mathbb{P}^{5}$ 中的这个点是 $V$ 在 Plücker 映射下的像.
\end{Example}

推广前例, 现在要描述以 Julius Plücker (1801--1868) 命名的 Plücker 映射
\[
  P_{k,n}:\mathrm{G}(k,n)\to\mathbb{P}^{\binom{n}{k}-1}.
\]
为了建立记号, 用大小为 $k$ 的子集 $I\subseteq\{1,\ldots,n\}$ 为 $\mathbb{P}^{\binom{n}{k}-1}$ 上的齐次坐标编号, 将相应的齐次坐标记为 $x_{I}$. 此外, 给定 $k\times n$ 矩阵 $M$ 和子集 $I\subseteq\{1,\ldots,n\}$, 令 $M_{I}$ 为 $M$ 中由 $I$ 编号的各列构成的 $k\times|I|$ 子矩阵. 对满足 $k\le n$ 的所有非负整数 $k,n$, 将 Plücker 映射定义如下.

\begin{Definition}{Plücker 映射}\leavevmode\label{def:plucker-maps}
\emph{Plücker 映射}
\[
  P_{k,n}:\mathrm{G}(k,n)\to\mathbb{P}^{\binom{n}{k}-1}
\]
是一个函数, 在 $V\in\mathrm{G}(k,n)$ 处的值如下给出: 选取 $V$ 的一组基, 用 $M(V)$ 表示以这些基元素为各行的 $k\times n$ 矩阵, 将 $P_{k,n}(V)\in\mathbb{P}^{\binom{n}{k}-1}$ 中由 $I$ 编号的齐次坐标的值设为 $\det(M(V)_{I})$.
\end{Definition}

这里应停下来确认 $P_{k,n}$ 是良定义的. 特别是, 应验证对 $V$ 的任意矩阵表示, $M(V)$ 的 $k\times k$ 子式中至少有一个非零, 并且点 $P_{k,n}(V)$ 与 $V$ 的基的选取无关 (Exercise~\ref{ex:plucker-map-well-defined}). 下一例说明这些观察.

\begin{Example}{计算 Plücker 映射}\leavevmode\label{ex:compute-plucker-map-for-plane-in-four-space}
\BookPairedNote{
令 $V\subseteq K^{4}$ 为 Example~\ref{ex:matrix-representations-of-same-plane} 的子空间, 将 $\mathbb{P}^{\binom{4}{2}-1}=\mathbb{P}^{5}$ 的坐标按字典序排列. 为简洁起见, 省去括号和逗号:
\[
  x_{12},\ x_{13},\ x_{14},\ x_{23},\ x_{24},\ x_{34}.
\]
}{
我们通常按字典序排列 $\mathbb{P}^{\binom{n}{k}-1}$ 的坐标, 就像在字典中排序一样, 只是用数字取代字母.
}
用基 $v_{1}=(2,1,-3,0)$ 和 $v_{2}=(1,0,1,1)$ 计算 $P_{2,4}(V)$, 得
\[
  [-1:5:2:1:1:-3]\in\mathbb{P}^{\binom{4}{2}-1},
\]
而用基 $w_{1}=(3,2,-7,-1)$ 和 $w_{2}=(3,1,-2,1)$ 得到 (同一个) 点
\[
  [-3:15:6:3:3:-9]\in\mathbb{P}^{\binom{4}{2}-1}.
\]
\end{Example}

与 Segre 映射的情形一样, 此时问 Plücker 映射是否正则还没有意义, 因为我们尚未赋予定义域射影簇结构. 事实上, 类似于 Definition~\ref{def:product-of-projective-varieties} 将积定义为射影簇的方式, 我们将通过证明 $P_{k,n}$ 是从 $\mathrm{G}(k,n)$ 到某个射影簇上的单射, 利用 Plücker 映射赋予 $\mathrm{G}(k,n)$ 射影簇结构. 先论证单射性, 再描述其像.

\begin{Proposition}{Plücker 映射是单射}\leavevmode\label{prop:plucker-maps-are-injective}
对满足 $0\le k\le n$ 的任意整数 $k,n$, Plücker 映射 $P_{k,n}$ 是单射.
\end{Proposition}
\begin{Proof}
设 $a\in\mathbb{P}^{\binom{n}{k}-1}$ 是 $P_{k,n}$ 的像中的一点. 我们论证存在唯一的线性子空间 $V\in\mathrm{G}(k,n)$, 使 $P_{k,n}(V)=a$. 选取非零坐标 $a_{I}\ne0$, 作缩放使 $a_{I}=1$. 若 $V\in\mathrm{G}(k,n)$ 映到 $a$, 则 $\det(M(V)_{I})\ne0$, 所以 $M(V)_{I}$ 可逆. 因此, 将 $M(V)$ 用 $M(V)_{I}^{-1}$ 相乘可知, 任一映到 $a$ 的子空间均有一个矩阵表示, 其由 $I$ 编号的各列组成一个单位矩阵. 现在论证, 映到 $a$ 的这样的矩阵表示只有一个.

设 $M$ 为某个 $V\in\mathrm{G}(k,n)$ 的矩阵表示, 使 $P_{k,n}(V)=a$ 且 $M_{I}$ 为单位矩阵. 令 $j\in\{1,\ldots,n\}\setminus I$; 我们要论证, 由 $P_{k,n}(V)=a$ 可唯一确定由 $j$ 编号的列 $M_{j}$. 对每个 $i\in I$, 考察集合 $I_{i}=(I\setminus\{i\})\cup\{j\}$. 由于 $M$ 映到 $a$, 有
\begin{equation}\label{eq:plucker-minor-determines-column-entry}
  a_{I_{i}}=\det(M_{I_{i}})\qquad\text{对所有 }i\in I.
\end{equation}
注意 $M_{I_{i}}$ 是从单位矩阵去掉第 $i$ 列, 再加入列 $M_{j}$ 得到的. 稍加思考可知, 这样一个矩阵的行列式无非是 $M_{j}$ 的第 $i$ 项 (相差一个符号). 特别地, 关系式 \eqref{eq:plucker-minor-determines-column-entry} 唯一确定了 $M_{j}$ 的各项. 因为 $j$ 是任意的, 所以由 $a$ 唯一确定 $M$, 证明完成.
\end{Proof}

下一步要描述 Plücker 映射下 $\mathrm{G}(k,n)$ 的像, 即 $\mathbb{P}^{\binom{n}{k}-1}$ 中的一个射影簇. 给定任一集合 $I\subseteq\{1,\ldots,n\}$ 和 $j\in\{1,\ldots,n\}$, 令 $I_{<j}$ 表示 $I$ 中小于 $j$ 的元素个数. 对每一对子集 $I,J\subseteq\{1,\ldots,n\}$, 若大小分别为 $k-1$ 和 $k+1$, 则定义相应的 \emph{Plücker 多项式}
\begin{equation}\label{eq:plucker-polynomial-by-two-index-sets}
  f_{I,J}=\sum_{j\in J}(-1)^{I_{<j}+J_{<j}}x_{I\cup\{j\}}x_{J\setminus\{j\}},
\end{equation}
其中约定当 $|I|\ne k$ 时 $x_{I}=0$ (特别地, 当 $j\in I$ 时, 和式的第 $j$ 项为零). 注意 $f_{I,J}$ 是关于 $\mathbb{P}^{\binom{n}{k}-1}$ 坐标的二次齐次多项式. 下面来看一个例子.

\begin{Example}{$\binom{4}{2}$ 个变量的 Plücker 多项式}\leavevmode\label{ex:plucker-polynomial-in-six-variables}
如 Example~\ref{ex:compute-plucker-map-for-plane-in-four-space} 所述, 将 $\mathbb{P}^{\binom{4}{2}-1}$ 的六个齐次坐标记为 $x_{12},x_{13},x_{14},x_{23},x_{24},x_{34}$. 若令 $I=\{1\}$ 且 $J=\{2,3,4\}$, 则得到 Plücker 多项式
\[
  f_{I,J}=-x_{12}x_{34}+x_{13}x_{24}-x_{14}x_{23}.
\]
另选任一对满足 $|I|=1$ 和 $|J|=3$ 的子集, 可以验证相应的 Plücker 多项式或者为零, 或者与上式相同, 至多相差一个符号. 因此, 在相差一个符号的意义下, $\binom{4}{2}$ 个变量的非零 Plücker 多项式只有一个.
\end{Example}

将 Example~\ref{ex:plucker-polynomial-in-six-variables} 的 Plücker 多项式 $f_{I,J}$ 在 Example~\ref{ex:compute-plucker-map-for-plane-in-four-space} 算出的点 $P_{k,n}(V)$ 处求值, 得
\[
  -(-1)(-3)+(5)(1)-(2)(1)=0.
\]
这是下面描述 Plücker 映射的像的结果的一个特例.

\begin{Proposition}{Plücker 映射与 Plücker 多项式}\leavevmode\label{prop:plucker-image-defined-by-plucker-polynomials}
对满足 $0\le k\le n$ 的任意整数 $k,n$, 有
\[
  P_{k,n}(\mathrm{G}(k,n))=
  \mathcal{V}\bigl(f_{I,J}\mid I,J\subseteq\{1,\ldots,n\}
  \text{ 且 }|I|=k-1,\ |J|=k+1\bigr).
\]
\end{Proposition}
\begin{Proof}
我们证明两个方向的包含关系.

($\subseteq$) 设 $V\in\mathrm{G}(k,n)$ 是 $K^{n}$ 的一个 $k$ 维线性子空间, 并令 $M$ 为 $V$ 的任一矩阵表示. 则对任意 $I,J$, 有
\[
  f_{I,J}(P_{k,n}(V))
  =\sum_{j\in J}(-1)^{I_{<j}+J_{<j}}
      \det(M_{I\cup\{j\}})\det(M_{J\setminus\{j\}}).
\]
令 $M_{I}^{i}$ 为从 $k\times(k-1)$ 矩阵 $M_{I}$ 删掉第 $i$ 行得到的矩阵, 令 $M_{J}^{i}$ 为将 $k\times(k+1)$ 矩阵 $M_{J}$ 的第 $i$ 行复制到最上方作为首行得到的矩阵. 将 $\det(M_{I\cup\{j\}})$ 按第 $j$ 列展开, 并将 $\det(M_{J}^{i})$ 按首行展开, 可验证 (Exercise~\ref{ex:determinant-expansions-for-plucker-equations})
\[
  \begin{aligned}
  &\sum_{j\in J}(-1)^{I_{<j}+J_{<j}}\det(M_{I\cup\{j\}})\det(M_{J\setminus\{j\}})\\
  &\qquad=\sum_{i=1}^{k}(-1)^{i}\det(M_{I}^{i})\det(M_{J}^{i}).
  \end{aligned}
\]
由于 $M_{J}^{i}$ 有两行相同, $\det(M_{J}^{i})=0$, 所以 $f_{I,J}$ 在 $P_{k,n}(V)$ 处消没.

($\supseteq$) 设对满足 $|I|=k-1$ 且 $|J|=k+1$ 的所有 $I,J\subseteq\{1,\ldots,n\}$, 有 $f_{I,J}(a)=0$. 固定一个使 $a_{I'}\ne0$ 的 $I'$, 将 $a$ 的齐次坐标缩放至 $a_{I'}=1$. 定义 $k\times n$ 矩阵 $M=(m_{i,j})$, 其中
\[
  m_{i,j}=(-1)^{I'_{<i}+I'_{<j}}a_{(I'\setminus\{i\})\cup\{j\}},
\]
若 $(I'\setminus\{i\})\cup\{j\}$ 的大小不是 $k$, 就令 $m_{i,j}=0$. 注意 $M_{I'}$ 就是单位矩阵, 因此 $M$ 的各行线性无关, 张成一个 $k$ 维线性子空间 $V\subseteq K^{n}$. 我们论证 $P_{k,n}(V)=a$.

为方便起见, 令 $P_{k,n}(V)=b$, 并缩放 $b$ 的齐次坐标, 使 $b_{I'}=a_{I'}=1$. 对任意 $i\in I'$ 及 $j\in\{1,\ldots,n\}\setminus I'$, 注意 $M_{(I'\setminus\{i\})\cup\{j\}}$ 与单位矩阵恰好只有一列不同. 计算得
\[
  \begin{aligned}
  b_{(I'\setminus\{i\})\cup\{j\}}
    &=\det(M_{(I'\setminus\{i\})\cup\{j\}})\\
    &=(-1)^{I'_{<i}+I'_{<j}}m_{i,j}=a_{(I'\setminus\{i\})\cup\{j\}},
  \end{aligned}
\]
其中第一个等式来自 $b$ 的定义, 第二个来自按第 $j$ 列展开行列式, 第三个来自 $m_{i,j}$ 的定义. 因而, 对满足 $|I\setminus I'|\le1$ 的任一子集 $I$, 有 $b_{I}=a_{I}$.

接下来对 $|I\setminus I'|$ 作归纳, 证明对所有 $I$, $b_{I}=a_{I}$. 上一段已证明基础情形. 为了建立归纳步骤, 固定 $\ell\ge2$, 假设只要 $|I\setminus I'|<\ell$ 就有 $b_{I}=a_{I}$; 现在固定某个满足 $|I\setminus I'|=\ell$ 的 $I$. 选取 $i\in I\setminus I'$, 考察 Plücker 多项式
\[
  f_{I\setminus\{i\},I'\cup\{i\}}
  =\sum_{j\in I'\cup\{i\}}(-1)^{(I\setminus\{i\})_{<j}+(I'\cup\{i\})_{<j}}
    x_{(I\setminus\{i\})\cup\{j\}}x_{(I'\cup\{i\})\setminus\{j\}}.
\]
此多项式在 $a$ 处消没 (由 $a$ 的定义), 并且也在 $b=P_{k,n}(V)$ 处消没 (由上面的 ($\subseteq$) 的证明). 将 $a$ 和 $b$ 代入这一关系, 由 $j=i$ 编号的项分别等于 $a_{I}$ 和 $b_{I}$ (相差一个符号), 其余每项都是由满足 $|J\setminus I'|<\ell$ 的子集 $J$ 编号的坐标的乘积. 由于对所有这样的 $J$ 有 $a_{J}=b_{J}$ (由归纳假设), 所以 $a_{I}=b_{I}$. 这就完成了归纳步骤, 从而完成证明.
\end{Proof}

Proposition~\ref{prop:plucker-image-defined-by-plucker-polynomials} 使我们能够将 Grassmann 簇实现为射影簇.

\begin{Definition}{作为射影簇的 Grassmann 簇}\leavevmode\label{def:grassmannian-as-projective-variety}
设 $k,n$ 为满足 $k\le n$ 的非负整数. \emph{Grassmann 簇} $\mathrm{G}(k,n)$ (作为射影簇) 是射影簇
\[
  P_{k,n}(\mathrm{G}(k,n))\subseteq\mathbb{P}^{\binom{n}{k}-1}.
\]
\end{Definition}

如本节开头所提到的, Grassmann 簇可用来研究射影空间中有关线性子簇的许多有趣的几何问题. 为了纪念对这些问题的早期研究作出许多贡献的 Hermann Schubert (1848--1911), 这一代数几何分支通常称为\emph{Schubert 演算}. 在下一例中, 我们回到本节开头提出的问题, 细节留作习题.

\begin{Example}{Schubert 演算初探}\leavevmode\label{ex:lines-meeting-four-fixed-projective-lines}
假设 $\mathbb{P}^{3}$ 中有四条固定直线 $L_{1},\ldots,L_{4}$, 想知道 $\mathbb{P}^{3}$ 中是否有第五条直线与每条 $L_{1},\ldots,L_{4}$ 相交. 由直线 $L\subseteq\mathbb{P}^{3}$ 对应于 $\mathrm{G}(2,4)$ 中的点, 可以定义四个子集 $X_{1},\ldots,X_{4}\subseteq\mathrm{G}(2,4)$:
\[
  X_{i}=\{L\in\mathrm{G}(2,4)\mid L\cap L_{i}\ne\varnothing\}.
\]
注意, 与每条 $L_{i}$ 相交的直线必在交集 $X_{1}\cap\cdots\cap X_{4}$ 中. 因此, 我们想知道 $X_{1}\cap\cdots\cap X_{4}$ 是否非空. 现在利用把 $\mathrm{G}(2,4)$ 解释为射影簇的观点:
\[
  \mathrm{G}(2,4)=\mathcal{V}(-x_{12}x_{34}+x_{13}x_{24}-x_{14}x_{23})
  \subseteq\mathbb{P}^{5}.
\]
对每个 $i=1,\ldots,4$, 可以验证 (Exercise~\ref{ex:lines-meeting-fixed-line-form-hyperplane-section}) 条件 $L_{i}\cap L\ne\varnothing$ 在 $L$ 的 Plücker 坐标上给出一个线性关系. 换句话说, 对每个 $i=1,\ldots,4$, 存在线性超平面 $H_{i}\subseteq\mathbb{P}^{5}$, 使
\[
  X_{i}=\mathrm{G}(2,4)\cap H_{i}.
\]
因此,
\[
  X_{1}\cap\cdots\cap X_{4}=\mathrm{G}(2,4)\cap(H_{1}\cap\cdots\cap H_{4}).
\]
$\mathbb{P}^{5}$ 中四个超平面的交给出一个至少一维的线性子簇 (Exercise~\ref{ex:codimension-of-linear-subvariety-intersection}). 由于任何这样的线性子簇必与 $\mathcal{V}(-x_{12}x_{34}+x_{13}x_{24}-x_{14}x_{23})$ 有非空交集 (Exercise~\ref{ex:positive-dimensional-linear-subvariety-meets-hypersurface}), 得 $X_{1}\cap\cdots\cap X_{4}\ne\varnothing$. 因此, $\mathbb{P}^{3}$ 中必有至少一条直线与任意给定的四条直线都相交.

事实上, 可以证明对 $L_{1},\ldots,L_{4}$ 的\textquotedbl{}几乎所有\textquotedbl{}选取, $\mathbb{P}^{3}$ 中恰有两条不同的直线与这四条直线都相交. 我们尚不具备严格论证所需的工具, 但基本想法是, 只要 $H_{1}\cap\cdots\cap H_{4}$ 是 $\mathbb{P}^{5}$ 中一条\textquotedbl{}一般的\textquotedbl{}直线, 由 $\mathrm{G}(2,4)$ 由单个二次多项式定义这一事实, 它与此直线的交恰是两个不同的点.
\end{Example}

本章已开始熟悉射影簇以及它们之间的映射. 此时很自然会像处理仿射簇那样, 转而讨论射影簇的维数和光滑性等性质. 然而, 鉴于我们已经在仿射情形为建立这些概念投入了不少工作, 如果能够利用这些工作, 而不必在这个不同的语境下重新发展同样的工具, 那当然再好不过. 为了建立将仿射工具引入射影情形所需的工具, 下一章将转向拟射影簇的概念, 它同时涵盖仿射和射影两种情形. 这个更一般的框架使我们能够把仿射簇视为射影簇的局部模型, 从而能将维数和光滑性的概念直接从仿射情形引入 (拟) 射影情形.

\BookExerciseHeading
\begin{BookExercises}
  \item\label{ex:linear-projective-subvarieties-correspond-to-vector-subspaces} 证明 $K^{n+1}$ 的 $(k+1)$ 维线性子空间的集合与 $\mathbb{P}^{n}$ 的 $k$ 维线性子簇的集合之间有保持包含关系的双射.
  \item\label{ex:characterize-linear-projective-subvarieties} 令 $L\subseteq\mathbb{P}^{n}$ 为一个子集. 证明下列条件等价.
  \begin{enumerate}[label=(\roman*)]
    \item $L$ 是 $\mathbb{P}^{n}$ 的 $k$ 维线性子簇.
    \item $L$ 是单射线性映射 $\mathbb{P}^{k}\to\mathbb{P}^{n}$ 的像.
    \item $L=\mathcal{V}(\ell_{1},\ldots,\ell_{n-k})$, 其中 $\ell_{1},\ldots,\ell_{n-k}\in K[x_{0},\ldots,x_{n}]$ 是线性无关的一次齐次多项式.
  \end{enumerate}
  \item\label{ex:codimension-of-linear-subvariety-intersection} 设 $L_{1},L_{2}$ 是 $\mathbb{P}^{n}$ 的线性子簇, 余维分别为 $d_{1},d_{2}$. 证明 $L_{1}\cap L_{2}$ 是余维至多为 $d_{1}+d_{2}$ 的线性子簇.
  \item\label{ex:grassmannian-of-hyperplanes-and-projective-space} 对任意 $n\ge0$, 描述 $\mathrm{G}(n,n+1)$ 与 $\mathbb{P}^{n}$ 之间的一个自然双射 (不使用 Plücker 映射).
  \item\label{ex:plucker-map-well-defined} 证明 Plücker 映射是良定义的. 特别地, 对任意 $k$ 维线性子空间 $V\subseteq K^{n}$, 证明以下结论:
  \begin{enumerate}
    \item 若 $M(V)$ 是 $V$ 的任一矩阵表示, 则对至少一个大小为 $k$ 的子集 $I\subseteq\{1,\ldots,n\}$, 有 $\det(M(V)_{I})\ne0$.
    \item 若 $M(V)$ 和 $M(V)'$ 是 $V$ 的两个矩阵表示, 则存在 $\lambda\in K\setminus\{0\}$, 使对每个大小为 $k$ 的子集 $I\subseteq\{1,\ldots,n\}$, 均有 $\det(M(V)_{I})=\lambda\det(M(V)'_{I})$.
  \end{enumerate}
  \item\label{ex:determinant-expansions-for-plucker-equations} 验证 Proposition~\ref{prop:plucker-image-defined-by-plucker-polynomials} 的证明中的第二条行间等式.
  \item\label{ex:positive-dimensional-linear-subvariety-meets-hypersurface} 令 $f\in K[x_{0},\ldots,x_{n}]$ 为齐次多项式, $L\subseteq\mathbb{P}^{n}$ 为正维数的线性子簇. 证明 $L\cap\mathcal{V}(f)\ne\varnothing$. 提示: 将 $L$ 中一条直线参数化, 再在这条直线上的点处对 $f$ 求值.
  \item\label{ex:lines-meeting-fixed-line-form-hyperplane-section} 对每个线性子空间 $V\in\mathrm{G}(2,4)$, 选取一个矩阵表示 $M(V)$. 对任意 $V_{1},V_{2}\in\mathrm{G}(2,4)$, 令 $M(V_{1},V_{2})$ 为将 $M(V_{1})$ 放在 $M(V_{2})$ 上方所得的 $4\times4$ 矩阵.
  \begin{enumerate}
    \item 证明 $\dim(V_{1}\cap V_{2})>0$ 当且仅当 $\det(M(V_{1},V_{2}))=0$.
    \item 证明 $4\times4$ 行列式 $\det(M(V_{1},V_{2}))$ 可表示为
    \[
      \det(M(V_{1},V_{2}))
      =\sum_{\substack{I\subseteq\{1,2,3,4\}\\|I|=2}}
       s(I)\det(M(V_{1})_{I})\det(M(V_{2})_{I^{c}}),
    \]
    其中 $I^{c}$ 表示补集 $\{1,2,3,4\}\setminus I$, 且
    \[
      s(I)=\begin{cases}
      1 & \text{若 }I=\{1,2\},\ \{1,4\},\ \{2,3\}\text{ 或 }\{3,4\},\\
      -1 & \text{若 }I=\{2,4\}\text{ 或 }\{3,4\}.
      \end{cases}
    \]
    \item 固定 $V_{0}\in\mathrm{G}(2,4)$, 并设 $X=\{V\in\mathrm{G}(2,4)\mid\dim(V\cap V_{0})>0\}$. 证明
    \[
      X=\mathrm{G}(2,4)\cap\mathcal{V}(\ell),
    \]
    其中 $\ell$ 是关于 $\mathbb{P}^{\binom{4}{2}-1}$ 的坐标的齐次线性多项式.
  \end{enumerate}
  将 $K^{4}$ 中的平面重新解释为 $\mathbb{P}^{3}$ 中的直线, 本题表明, $\mathbb{P}^{3}$ 中与一条固定直线相交的直线的集合, 是 $\mathrm{G}(2,4)$ 与一个超平面的交.
  \item\label{ex:grassmannian-affine-charts-and-irreducibility} 通过 Plücker 映射将 $\mathrm{G}(k,n)$ 视为一个射影簇, 并令 $I$ 为 $\{1,\ldots,n\}$ 的大小为 $k$ 的一个子集.
  \begin{enumerate}
    \item 证明仿射限制
    \[
      \mathrm{G}(k,n)\setminus\mathcal{V}(x_{I})
        \subseteq\mathbb{A}^{\binom{n}{k}-1}
    \]
    作为仿射簇与 $\mathbb{A}^{k(n-k)}$ 同构.
    \item 证明 $\mathrm{G}(k,n)\setminus\mathcal{V}(x_{I})$ 的射影闭包是 $\mathrm{G}(k,n)$.
    提示: 给定任意 $V\in\mathcal{V}(x_{I})$, 构造函数 $f:\mathbb{A}^{1}\to\mathrm{G}(k,n)$, 使得 $f(0)=V$, 且当 $a\ne0$ 时 $f(a)\notin\mathcal{V}(x_{I})$. 简化行阶梯形或许有用. 然后论证, 任何对 $a\ne0$ 都在 $f(a)$ 处消没的多项式也必在 $f(0)$ 处消没.
    \item 用 (a) 和 (b) 论证 $\mathrm{G}(k,n)$ 不可约.
  \end{enumerate}
  在正式定义射影簇的维数后, 我们还可利用本题证明 $\mathrm{G}(k,n)$ 的维数为 $k(n-k)$.
\end{BookExercises}
